\section{Related work}
\label{sec:related}

\paragraph{Input format.} Sui et al.~\cite{sui2024tablemeetsllm} compare
table serializations (Markdown, HTML, JSON, XML, \ldots) on structural
understanding tasks and find that HTML with format explanations does well.
He et al.~\cite{he2024promptformatting} show that the same prompt
written as plain text, Markdown, JSON or YAML can move GPT accuracy by tens
of points. Sclar et al.~\cite{sclar2024formatting} find similar sensitivity
to formatting choices that carry no meaning. These studies vary the
\emph{prompt}. We vary the \emph{document}: what the model is asked
\emph{about}, in the formats it comes in from files.

\paragraph{Output format.} Tam et al.~\cite{tam2024speakfreely} show
that forcing JSON or XML output can hurt reasoning. FOFO~\cite{xia2024fofo}
measures whether models follow domain-specific formats. Both score the
\emph{text}. Neither builds the file nor asks what a valid DOCX, PDF, PPTX
or XLSX costs, when the model can only produce it through code, markup or
raw OOXML.

\paragraph{Document conversion and context.} Pipelines turn PDFs and
office files into model input with extractors and layout models
(Docling~\cite{auer2024docling}). The choice is usually made for
extraction quality, not for what the result costs the model. Token cost
differs by language under the same tokenizer~\cite{petrov2023tokenizer},
and models use long contexts unevenly~\cite{liu2024lost}. We measure the
analogous cost across formats, and a fixed-budget condition in which markup
competes with content for the context.
